
\PuzzleChapterIntro{The Ledger of Five Seals}{Mixed \textbullet\ Theorems \textbullet\ Truth Vector}{Five independent claims are sealed into a single lock: only the correct truth values form the five-bit signature. Each statement is self-contained, but the door opens only if you classify all five correctly.}

\Lore{
The Ledger of the Five Seals is written in a cramped hand, the kind of handwriting that only appears when someone is trying to make a rule survive an emergency.

It lists five charms. Each charm can be set to one of two states.
The margin note—underlined twice—says:

\begin{quote}
\textit{Only the full five-bit signature opens the door. Nothing else.}
\end{quote}

The clerk asks for the signature itself and nothing more.
}

\Problem

For each of the following five statements, assign a truth value: assign $1$ if the statement is true and $0$ if it is false.
Express your final answer as a $1\times 5$ matrix
\[
T=[t_A,t_B,t_C,t_D,t_E].
\]

\begin{enumerate}[label=\textbf{\Alph*.},leftmargin=*]

\item \textbf{No-new-prime claim for odd prime exponents.}

Let $m$ and $n$ be coprime positive \emph{odd} integers with $m>n$, and let $k\ge 5$ be an \emph{odd prime}.
\emph{Claim:} For every prime $r$ dividing $m^k-n^k$, there exists a \emph{proper divisor} $j$ of $k$ (so $1\le j<k$ and $j\mid k$) such that
\[
r\mid m^j-n^j.
\]

\item \textbf{Cubic non-residue inverse-sum.}

Let $p$ be a prime with $p\equiv 1\pmod 3$ and $p>7$. Define
\[
CR_p=\{a\in\{1,2,\dots,p-1\}:\ a\text{ is \emph{not} a cubic residue modulo }p\}.
\]
(For $a\in\{1,\dots,p-1\}$, write $a^{-1}$ for the multiplicative inverse of $a$ modulo $p$.)
\emph{Claim:}
\[
\sum_{a\in CR_p} a^{-1} \equiv 0 \pmod p.
\]

\item \textbf{Triangle-free graph edge bound.}

Let $G=(V,E)$ be a simple graph on $n$ vertices.
\emph{Claim:} If $G$ is triangle-free, then
\[
|E|\le \left\lfloor \frac{n^2}{4}\right\rfloor.
\]

\item \textbf{Subset pairing symmetry lemma (interior $k$ only).}

Let $n\ge 3$ and let $S$ be a finite set with $|S|=n$, and let $f:S\to\mathbb{R}$.
For each $k$ with $0\le k\le n$, define
\[
F(k)=\sum_{\substack{T\subseteq S\\ |T|=k}}\ \sum_{x\in T} f(x).
\]
\emph{Claim:}
\[
\bigl(\forall\,k\in\{1,2,\dots,n-1\},\ F(k)=F(n-k)\bigr)\ \Longleftrightarrow\ \left(\sum_{x\in S} f(x)=0\right).
\]

\item \textbf{Ultimate monotonic subsequence lemma.}

Let $n$ be a positive integer and let $A=(a_1,a_2,\dots,a_{n^2+1})$ be any sequence of $n^2+1$ distinct real numbers.
Let $\ell_{\mathrm{inc}}(A)$ and $\ell_{\mathrm{dec}}(A)$ denote the lengths of the longest strictly increasing and strictly decreasing subsequences, respectively.
\emph{Claim:}
\[
\max\{\ell_{\mathrm{inc}}(A),\ell_{\mathrm{dec}}(A)\}\ge n+1.
\]

\end{enumerate}

\VaultDirective{Output the 5-bit signature $t_A t_B t_C t_D t_E$.}

\SolutionPage
\small
\begin{enumerate}[label=\textbf{\Alph*.},leftmargin=*]

\item \textbf{False.}
Since $k$ is an odd prime, its only proper divisor is $j=1$, so the claim reduces to:

\emph{every prime divisor $r$ of $m^k-n^k$ must divide $m-n$.}

Counterexample: take $m=3$, $n=1$, $k=5$ (coprime, odd, and $k$ is an odd prime).
Then
\[
3^5-1^5=243-1=242=2\cdot 11^2.
\]
The prime $11$ divides $3^5-1^5$, but $11\nmid (3-1)=2$.
So the claim fails.

\item \textbf{True.}
Let $G=\mathbb{F}_p^\times$, a cyclic group of order $p-1$.
Since $p\equiv 1\pmod 3$, the set of nonzero cubes forms a subgroup $H\le G$ of index $3$, with
\[
|H|=\frac{p-1}{3}>1 \quad (\text{since }p>7).
\]
Thus $CR_p$ is exactly the union of the two nontrivial cosets of $H$.

First, $\sum_{h\in H} h\equiv 0\pmod p$:
the elements of $H$ are precisely the roots in $\mathbb{F}_p$ of $x^{|H|}-1$, whose $x^{|H|-1}$ coefficient is $0$, so the sum of its roots is $0$.
Hence for any $g\in G$, the coset $gH$ sums to
\[
\sum_{x\in gH} x = g\sum_{h\in H} h \equiv 0 \pmod p.
\]
Therefore the sum of all elements of $CR_p$ is $0\pmod p$.

Now observe that inversion permutes $CR_p$:
if $a\notin H$, then $a\in gH$ or $a\in g^2H$ for some $g\notin H$, and $a^{-1}\in (gH)^{-1}=g^{-1}H=g^2H$, still outside $H$.
So $\{a^{-1}:a\in CR_p\}=CR_p$, and
\[
\sum_{a\in CR_p} a^{-1}\equiv \sum_{a\in CR_p} a \equiv 0\pmod p.
\]

\item \textbf{True.}
This is Mantel’s theorem: among triangle-free graphs on $n$ vertices, the maximum number of edges is attained by the complete bipartite graph with parts as equal as possible, giving $\lfloor n^2/4\rfloor$.

\item \textbf{True.}
For $k\ge 1$, each $x\in S$ appears in exactly $\binom{n-1}{k-1}$ subsets $T\subseteq S$ of size $k$, so
\[
F(k)=\binom{n-1}{k-1}\sum_{x\in S} f(x).
\]
In particular,
\[
F(1)=\sum_{x\in S} f(x),
\qquad
F(n-1)=\binom{n-1}{n-2}\sum_{x\in S} f(x)=(n-1)\sum_{x\in S} f(x).
\]
If $F(1)=F(n-1)$ and $n\ge 3$, then
\[
\sum f(x)=(n-1)\sum f(x)\ \Longrightarrow\ (n-2)\sum f(x)=0\ \Longrightarrow\ \sum f(x)=0.
\]
Conversely, if $\sum f(x)=0$, then $F(k)=0$ for every $k\ge 1$, so $F(k)=F(n-k)$ holds for all $k=1,\dots,n-1$.

\item \textbf{True.}
This is exactly the Erd\H{o}s--Szekeres monotone subsequence theorem for sequences of length $n^2+1$ of distinct real numbers.

\end{enumerate}

Therefore,
\[
T=[0,1,1,1,1].
\]

\FinalAnswer{01111}

\NextPage
